\documentclass[11pt,letterpaper]{article}
\usepackage[margin=0.98in]{geometry}
\usepackage[T1]{fontenc}
\usepackage{lmodern}
\usepackage{microtype}
\usepackage{amsmath,amssymb,amsthm,mathtools}
\usepackage{enumitem}
\usepackage{hyperref}
\hypersetup{colorlinks=true, linkcolor=blue!45!black, citecolor=blue!45!black, urlcolor=blue!45!black, pdftitle={Sharp harmonic cycle-length inequalities for complete multipartite graphs}}
\usepackage{xcolor}
\usepackage{fancyhdr}
\usepackage{lastpage}
\usepackage{booktabs}
\setlength{\headheight}{14pt}
\pagestyle{fancy}
\fancyhf{}
\fancyhead[L]{\footnotesize Harmonic cycle-length inequalities}
\fancyhead[R]{\footnotesize Preprint --- 9 October 2026}
\fancyfoot[C]{\footnotesize \thepage\,/\,\pageref{LastPage}}
\setlength{\parindent}{1.2em}
\setlength{\parskip}{0.14em}
\newtheorem{theorem}{Theorem}[section]
\newtheorem{lemma}[theorem]{Lemma}
\newtheorem{proposition}[theorem]{Proposition}
\newtheorem{corollary}[theorem]{Corollary}
\newtheorem{remark}[theorem]{Remark}
\newcommand{\cset}{\mathcal{C}}
\newcommand{\hc}{\mathsf{s}}
\newcommand{\HH}{H}
\newcommand{\edges}{e}
\newcommand{\vcover}{\tau}
\newcommand{\N}{\mathbb N}
\title{\bfseries Sharp Harmonic Cycle-Length Inequalities\\for Complete Multipartite Graphs\\[5pt]
\large With bounded-cover and near-balanced extensions}
\author{\textit{Author name and affiliation to be inserted}}
\date{9 October 2026}
\begin{document}
\maketitle
\begin{abstract}
For a finite simple graph $G$, let $\cset(G)$ be the set of distinct lengths of its simple cycles and let
$\hc(G)=\sum_{\ell\in\cset(G)}\ell^{-1}$.  A conjecture of Erd\H{o}s asserts that if $n\ge 2d$ and
$\edges(G)\ge d(n-d)$, then $\hc(G)\ge \frac12(H_d-1)$, with equality only for $K_{d,n-d}$.
We prove a sharp version for every complete multipartite graph, for every integer $d\ge 2$,
under the weaker hypothesis $\edges(G)>(d-1)(n-d+1)$.  The proof is based on an exact
classification of the cycle-length spectrum of complete multipartite graphs.
We also establish the original inequality, with its equality characterization, for graphs
having a vertex cover of size at most $d$ under an explicit order hypothesis, and for
arbitrary graphs in the near-balanced range $d\ge4$, $2d\le n\le2d+4$.
The latter proof invokes classical theorems of Brandt, Erd\H{o}s--Gallai, and Moon--Moser.
These are restricted-family results, not a proof of the unrestricted conjecture.
\end{abstract}
\noindent\textbf{Keywords.} Cycle lengths; harmonic sums; complete multipartite graphs;
Erd\H{o}s--Hajnal conjecture; extremal graph theory; pancyclicity.\\
\noindent\textbf{2020 Mathematics Subject Classification.} 05C38, 05C35.

\section{Introduction}
We work throughout with finite, undirected, simple graphs.  For a graph $G$ of order $n$, define
\[
\cset(G)=\{\ell\in\{3,\dots,n\}:G\text{ contains a cycle of length }\ell\},
\qquad \hc(G)=\sum_{\ell\in\cset(G)}\frac1\ell.
\]
Write $H_t=\sum_{j=1}^{t}j^{-1}$ for the $t$th harmonic number and
\[
B_d=\sum_{j=2}^{d}\frac1{2j}=\frac12(H_d-1).
\]
The complete bipartite graph $K_{d,n-d}$ has
$\cset(K_{d,n-d})=\{4,6,\ldots,2d\}$ whenever $n\ge 2d$; hence
$\hc(K_{d,n-d})=B_d$.

A sharp form of the Erd\H{o}s--Hajnal harmonic cycle-length problem asks whether
\begin{equation}\label{eq:erdos}
 n\ge 2d,\quad \edges(G)\ge d(n-d)
 \qquad\Longrightarrow\qquad
 \hc(G)\ge B_d,
\end{equation}
with equality only when $G\cong K_{d,n-d}$; see \cite{MMPS,EP}.
Montgomery, Milojevi\'c, Pokrovskiy, and Sudakov \cite{MMPS} recently established the
stronger edge-threshold version for all sufficiently large $d$.  It is important that this
known theorem does not state a uniform result for every $d\ge2$.

Our main result proves their sharp edge-threshold inequality for 
\emph{all} $d\ge2$ under the complete multipartite hypothesis.  We give an exact
cycle-spectrum formula rather than assuming a general pancyclicity theorem.
Two additional results concern graphs whose structure is less rigid: one has an
unbounded number of vertices but a small vertex cover; the other applies to
arbitrary graphs in a short near-balanced interval of orders.

\begin{theorem}[Sharp complete multipartite case]\label{thm:main}
Let $d\ge 2$ and $n\ge2d$ be integers.  If $G$ is a complete multipartite graph
on $n$ vertices and
\begin{equation}\label{eq:sharp}
\edges(G)>(d-1)(n-d+1),
\end{equation}
then
\[
\hc(G)\ge B_d.
\]
Equality holds if and only if $G\cong K_{d,n-d}$.
The strict threshold in \eqref{eq:sharp} cannot be weakened to a non-strict one,
even within complete multipartite graphs.
\end{theorem}

The two extensions are as follows.
\begin{theorem}[Bounded vertex cover]\label{thm:cover}
Let $d\ge2$ and $n\ge2d+\binom d2$.  If an $n$-vertex graph $G$ satisfies
$\vcover(G)\le d$ and $\edges(G)\ge d(n-d)$, then
$\hc(G)\ge B_d$, with equality if and only if $G\cong K_{d,n-d}$.
\end{theorem}
Here $\vcover(G)$ denotes the minimum size of a vertex cover.

\begin{theorem}[Arbitrary graphs in a near-balanced range]\label{thm:band}
If $d\ge4$ and $2d\le n\le2d+4$, then every $n$-vertex graph $G$ with
$\edges(G)\ge d(n-d)$ satisfies $\hc(G)\ge B_d$, with equality
if and only if $G\cong K_{d,n-d}$.
\end{theorem}

\section{The complete multipartite cycle spectrum}
We start with a precise classification that supplies the principal combinatorial
input.  The graph $G$ is \emph{complete multipartite} if its vertex set is a
partition into nonempty independent classes, and two vertices are adjacent precisely
when they belong to different classes.

\begin{lemma}[Capped selection]\label{lem:capacity}
Let $a_1,\ldots,a_r$ be positive integers, $r\ge3$, with sum $n$, and put
$b=\max_i a_i$.  Suppose $b<n/2$.  Then for every $3\le\ell\le n$, with
$q=\lfloor\ell/2\rfloor$, one has
\begin{equation}\label{eq:capacity}
\sum_{i=1}^{r}\min(a_i,q)\ge\ell.
\end{equation}
\end{lemma}
\begin{proof}
If all $a_i\le q$, the left side equals $n\ge\ell$.
If exactly one class has size greater than $q$, that class contributes $q$.
The other classes together have $n-b>b>q$ vertices, so they contribute
at least $q+1$; here their sizes are individually at most $q$.
Thus the sum is at least $2q+1\ge\ell$.
Finally, if two or more classes exceed $q$, their combined contribution
is at least $2q$.  This is enough when $\ell=2q$.  When $\ell=2q+1$,
there is at least one further nonempty class, unless three or more already
exceed $q$, in either event contributing at least one extra unit.
This proves \eqref{eq:capacity}.
\end{proof}

\begin{proposition}[Exact cycle-length spectrum]\label{prop:spectrum}
Let $G$ be complete multipartite, with $r\ge2$ nonempty classes,
$n$ vertices, largest class size $b$, and $s=n-b$.
\begin{enumerate}[label=(\roman*),itemsep=0.2em]
\item If $r=2$, then $\cset(G)=\{4,6,\ldots,2s\}$.
\item If $r\ge3$, then
\begin{equation}\label{eq:spectrum}
\cset(G)=\{3,4,\ldots,\min(n,2s)\}.
\end{equation}
\end{enumerate}
\end{proposition}
\begin{proof}
For (i), $G=K_{b,s}$ with $b\ge s$, so the alternating cycles have
exactly the lengths displayed, and a bipartite graph has no odd cycle.

For (ii), first suppose $b\ge s$.  There are at least two classes
outside the largest class, so $s\ge2$.
For $2\le q\le s$ select $q$ vertices in the largest class and
$q$ vertices outside it.  Their complete bipartite cross-edges
contain a $2q$-cycle.
For $1\le q\le s-1$, select $q$ vertices in the largest class and
$q+1$ outside vertices in such a way that no outside class
contributes more than $q$ vertices.  This is possible because the
outside classes have total size $s\ge q+1$ and at least two are
nonempty: if one exceeds $q$, the others contribute at least one;
otherwise all $s$ vertices are eligible.  The selected $(2q+1)$-vertex
induced graph has every class of size at most $q$ and hence minimum
degree at least $(2q+1)-q=q+1$.  Dirac's Hamiltonicity theorem
\cite{Dirac} (also valid for order three) supplies a cycle of length
$2q+1$.  All lengths $3,4,\ldots,2s$ occur.

No cycle can use more than $s$ vertices from the largest class:
vertices of that class are independent, and successive such vertices
around a cycle must be separated by distinct vertices outside the class.
Therefore every cycle has length at most $2s$.  As $b\ge s$,
$\min(n,2s)=2s$, proving \eqref{eq:spectrum} in this case.

Now suppose $b<s$, or equivalently $b<n/2$.  Fix $3\le\ell\le n$
and let $q=\lfloor\ell/2\rfloor$.  Lemma~\ref{lem:capacity}
allows the selection of exactly $\ell$ vertices, with at most $q$
from any class.  The induced complete multipartite graph has minimum
degree at least $\ell-q=\lceil\ell/2\rceil$.  Dirac's theorem
therefore supplies an $\ell$-cycle.  Every length from 3 through
$n$ occurs, completing the proof.
\end{proof}

\begin{corollary}[Exact harmonic sum]\label{cor:formula}
With the notation of Proposition~\ref{prop:spectrum},
\[
\hc(G)=
\begin{cases}
\frac12(H_s-1),&r=2,\\[3pt]
H_{\min(n,2s)}-\frac32,&r\ge3.
\end{cases}
\]
\end{corollary}
\begin{proof}
Sum the reciprocals of the lengths listed in
Proposition~\ref{prop:spectrum}.
\end{proof}

\section{Proof of the sharp multipartite theorem}
\begin{proof}[Proof of Theorem~\ref{thm:main}]
Let the class sizes of $G$ be $a_1\ge\cdots\ge a_r>0$, where
$r\ge2$, $b=a_1$, and $s=n-b$.  The case $r=1$ is ruled out by
\eqref{eq:sharp}.

Suppose first that $r=2$.  The graph is $K_{s,n-s}$, with
$s\le n/2$ and $\edges(G)=s(n-s)$.  The function
$x\mapsto x(n-x)$ is strictly increasing on integers
$0\le x\le n/2$.
Because
\[
(s(n-s))\ >\ (d-1)(n-d+1),
\]
we must have $s\ge d$.  Hence Corollary~\ref{cor:formula} gives
\[
\hc(G)=\frac12(H_s-1)\ge\frac12(H_d-1)=B_d,
\]
with equality exactly when $s=d$, i.e., $G\cong K_{d,n-d}$.

Suppose now that $r\ge3$.  Every edge has at least one endpoint
outside the largest class, so $\edges(G)\le sn$.
We claim that $2s\ge d$.  If not, then $s\le(d-1)/2$, yielding
\[
\edges(G)\le sn\le\frac{(d-1)n}{2}
 < (d-1)(n-d+1),
\]
because $n\ge2d$ implies $n-d+1>n/2$.
This contradicts \eqref{eq:sharp}.  Therefore $2s\ge d$.
Set $L=\min(n,2s)$.  Since $n\ge2d$, it follows that $L\ge d$.
Moreover $s\ge2$, hence $L\ge4$ and $H_L\ge H_4=25/12>2$.
Using Corollary~\ref{cor:formula},
\[
\hc(G)=H_L-\frac32
>\frac12(H_L-1)\ge\frac12(H_d-1)=B_d.
\]
Thus equality cannot occur when $r\ge3$.

Finally, the threshold is sharp: $K_{d-1,n-d+1}$ has
$(d-1)(n-d+1)$ edges and
\[
\hc(K_{d-1,n-d+1})=B_{d-1}<B_d,
\]
where $B_1=0$ for $d=2$.  This completes the proof.
\end{proof}

\begin{corollary}[The conjectured edge threshold inside the class]\label{cor:target}
If $G$ is complete multipartite, $n\ge2d$, and
$\edges(G)\ge d(n-d)$, then $\hc(G)\ge B_d$,
with equality only for $K_{d,n-d}$.
\end{corollary}
\begin{proof}
For $n\ge2d$,
\[
d(n-d)-(d-1)(n-d+1)=n-2d+1>0.
\]
Apply Theorem~\ref{thm:main}.
\end{proof}

\section{A bounded vertex-cover extension}
\begin{proof}[Proof of Theorem~\ref{thm:cover}]
Let $S$ be a vertex cover of minimum size $q\le d$ and put
$T=V(G)\setminus S$.  The set $T$ is independent, whence
\begin{equation}\label{eq:cover-upper}
\edges(G)\le q(n-q)+\binom q2
=qn-\frac{q(q+1)}2.
\end{equation}
The right side is increasing for integral $0\le q\le d$ under
$n\ge2d$.  If $q\le d-1$, then
\[
\edges(G)\le(d-1)n-\frac{d(d-1)}2< d(n-d)
\]
where the strict inequality follows from
$n\ge2d+\binom d2>\frac{d(d+1)}2$.
This contradicts our edge hypothesis, so $q=d$.

Write $Y=V(G)\setminus S$, and let $r=\edges(G[S])$.
Let $t$ count the edges missing between $S$ and $Y$.
The two independent- and cross-edge contributions give
\[
\edges(G)=d(n-d)-t+r\ge d(n-d),
\quad\text{so}\quad 0\le t\le r\le\binom d2.
\]
At most $t$ vertices of $Y$ are incident with a missing
cross-edge.  Consequently at least
\[
(n-d)-t\ge n-d-\binom d2\ge d
\]
vertices of $Y$ are adjacent to every vertex of $S$.
These vertices and $S$ contain $K_{d,d}$, which contributes
every even cycle length $4,6,\ldots,2d$.  Hence $\hc(G)\ge B_d$.

If $r>0$, an edge inside $S$, together with any vertex of $Y$
adjacent to all of $S$, forms a triangle.  This gives the
additional length three and strictly increases $\hc(G)$.
If $r=0$, then $t=0$; as $Y$ is independent, the graph is
exactly $K_{d,n-d}$.  The converse equality is immediate.
\end{proof}

\begin{remark}
The sufficient size condition in Theorem~\ref{thm:cover} is not
asserted optimal.  Its conclusion applies to graphs with arbitrary
edges inside $S$ and nonuniform neighborhoods in $Y$, so it does
not follow merely from Theorem~\ref{thm:main}.
\end{remark}

\section{A near-balanced theorem for arbitrary graphs}
We give the external results used in this section in the precise
forms needed below.  Brandt's theorem \cite{Brandt}, also recorded
in \cite[Theorem 2.5]{MMPS}, states that a nonbipartite $n$-vertex
graph with more than $(n-1)^2/4+1$ edges has cycles of every length
from 3 to its circumference.  Erd\H{o}s--Gallai \cite{EG}
shows that a graph with $m$ edges and $n$ vertices has a cycle
of length at least $2m/(n-1)$ whenever $m\ge n$.
The Moon--Moser bipartite Ore condition \cite[Corollary 3]{MoonMoser}
states that a balanced bipartite graph with $q$ vertices in each class
is Hamiltonian if every nonadjacent pair $x,y$ in opposite classes
satisfies $\deg(x)+\deg(y)>q$.

\begin{lemma}[Robust balanced bipartite Hamiltonicity]\label{lem:hami}
Let $q\ge3$.  If at most $q-2$ edges are removed from $K_{q,q}$,
then the remaining graph has a Hamiltonian cycle.
\end{lemma}
\begin{proof}
Write $t\le q-2$ for the number of missing edges.  For
nonadjacent opposite-class vertices $x,y$, let $m_x,m_y$ denote
the numbers of their missing cross-edges.  The absent edge $xy$
is counted twice in $m_x+m_y$, whereas every other missing edge
is counted at most once.  Thus $m_x+m_y\le t+1$, whence
\[
\deg(x)+\deg(y)=2q-(m_x+m_y)
\ge 2q-t-1\ge q+1>q.
\]
The cited Moon--Moser criterion gives a Hamiltonian cycle.
\end{proof}

\begin{lemma}[Covering four missing cross-edges]\label{lem:four}
The edge set of any bipartite graph having at most four edges
can be covered by at most two vertices from each bipartition class.
\end{lemma}
\begin{proof}
If some vertex has degree at least three, select that vertex;
at most one edge remains uncovered, and one endpoint of that edge
completes the cover.  If the maximum degree is two, select a
vertex of degree two.  Of the at most two remaining edges, choose
an endpoint on the side opposite the selected vertex for each,
using one selected vertex on the first side and at most two on
the other.  Finally, if every vertex has degree at most one, the
edges form a matching of size at most four.  Choose an endpoint
on the first side of at most two matching edges and an endpoint
on the second side of the remaining edges.
\end{proof}

\begin{proof}[Proof of Theorem~\ref{thm:band}]
Write $n=2d+r$ with $0\le r\le4$.

\emph{Case 1: $G$ is nonbipartite.}
A direct calculation yields
\begin{equation}\label{eq:brandt-diff}
 d(n-d)-\left(\frac{(n-1)^2}{4}+1\right)
 =\frac{4d-4-(r-1)^2}{4}>0,
\end{equation}
since $d\ge4$ and $(r-1)^2\le9$.
Brandt's criterion therefore supplies cycles of all lengths
from 3 to the circumference $c(G)$.
By the Erd\H{o}s--Gallai estimate,
\[
 c(G)\ge\frac{2\edges(G)}{n-1}
 \ge\frac{2d(n-d)}{n}\ge d.
\]
(The last inequality uses $n\ge2d$.)  Hence
\[
\hc(G)\ge H_d-\frac32
>\frac12(H_d-1)=B_d,
\]
because $H_d\ge H_4=25/12>2$.

\emph{Case 2: $G$ is bipartite.}
Let its bipartition class sizes satisfy $a\le b$, $a+b=n$.
Since $\edges(G)\le ab$ and $\edges(G)\ge d(n-d)$,
the monotonicity of $a(n-a)$ for $a\le n/2$ implies $a\ge d$.
Thus
\[
 a=d+j,\quad b=d+r-j,\quad
 0\le j\le\lfloor r/2\rfloor.
\]
The number $t$ of missing cross-edges satisfies
\begin{equation}\label{eq:missing}
 t=ab-\edges(G)\le j(r-j).
\end{equation}
If $j=0$, then $t=0$ and $G=K_{d,n-d}$.
Otherwise $(r,j)$ belongs to
\[
\{(2,1),(3,1),(4,1),(4,2)\},
\]
so $t\le1,2,3,4$, respectively.

We first exhibit a $K_{d,d}$ subgraph.  If $j=1$, delete
from the larger class an endpoint of each missing edge;
this uses at most $t\le r-1$ vertices, leaving at least
$(d+r-1)-(r-1)=d$ there and all $d+1$ vertices on
the other side.  Choose $d$ in each class.
If $(r,j)=(4,2)$, both original bipartition classes have
size $d+2$.  Apply Lemma~\ref{lem:four} to the graph of
missing cross-edges, delete its covering vertices (at most
two from each class), and choose $d$ surviving vertices
from each class.  Thus $K_{d,d}\subseteq G$ in every case
with $j>0$, and consequently
\[
\{4,6,\ldots,2d\}\subseteq\cset(G).
\]

It remains to find an additional cycle length.  For $j=1$,
choose $d+1$ vertices from each class.  The induced balanced
bipartite graph is obtained from $K_{d+1,d+1}$ by removing
at most $t\le r-1\le3$ edges.  Set $q=d+1\ge5$;
then $t\le3\le q-2$, so Lemma~\ref{lem:hami}
produces a cycle of length $2d+2$.
For $(r,j)=(4,2)$, the entire graph has classes of size
$q=d+2\ge6$ and at most four missing edges.  As
$4\le q-2$, Lemma~\ref{lem:hami} supplies a cycle of
length $2d+4$.
Both lengths exceed $2d$; the harmonic sum is therefore
strictly larger than $B_d$.  The sole equality case is
$j=0$, namely $K_{d,n-d}$.
\end{proof}

\section{A reduction and the remaining regime}
The elementary observation below is useful when searching
for a minimal counterexample to \eqref{eq:erdos} at fixed $d$.
\begin{proposition}[Deletion of a low-degree vertex]\label{prop:del}
Fix $d\ge2$.  Suppose $n>2d$, $G$ has $n$ vertices,
$\edges(G)\ge d(n-d)$, and $G$ has a vertex $v$ with
$\deg_G(v)\le d$.  Then $G-v$ has $n-1\ge2d$ vertices,
$\edges(G-v)\ge d((n-1)-d)$, and $\hc(G-v)\le\hc(G)$.
\end{proposition}
\begin{proof}
Removing $v$ deletes at most $d$ edges and cannot introduce
new cycle lengths.
\end{proof}
Thus a minimum-order counterexample to the inequality
(at a fixed $d$, if one exists) must either have order $2d$
or have minimum degree at least $d+1$.

The results here do not cover arbitrary graphs of unbounded order
for each fixed small $d$.  In particular, neither their correctness
nor the exact finite checks below constitutes a solution of the
unrestricted Erd\H{o}s problem.  The complete multipartite cycle
spectrum is elementary; whether the weighted consequence already
appears implicitly in published literature requires additional
specialist review.  We therefore make no priority claim.

\section*{Computational verification (limited scope)}
The supplemental program \texttt{verify\_graph\_atlas.py} enumerates
all isomorphism types of simple graphs on at most seven vertices
using NetworkX's Graph Atlas.  It computes the cycle-length spectrum
and reciprocal sum with exact rational arithmetic, then compares
against $B_d$ for all admissible $2\le d\le n/2$ and edge thresholds.
A local execution found 1,194 qualifying graph--parameter tests
and no violations; the only equality isomorphism type in each
parameter cell is $K_{d,n-d}$.  These finite tests do not directly
cover Theorem~\ref{thm:band}, which starts at $n=8$, and do not
replace its proof.

\section*{Research status and attribution}
This manuscript establishes precisely the restricted theorems
stated above, conditional only on the cited standard theorems
of graph theory.  It does \emph{not} settle Erd\H{o}s Problem \#65.
The literature search for this draft was targeted rather than
exhaustive; the claimed special-case statements have not been
independently certified as new or peer-reviewed.  The author
placeholder must be replaced by accurate authorship and affiliation
information before public release.

\begin{thebibliography}{9}
\bibitem{MMPS}
R. Montgomery, A. Milojevi\'c, A. Pokrovskiy, and B. Sudakov,
\emph{Minimising the harmonic sum of cycle lengths},
arXiv:2609.26401 (2026).
\url{https://arxiv.org/abs/2609.26401}.
\bibitem{EP}
\emph{Erd\H{o}s Problems}, Problem \#65, ``Harmonic sums of cycle lengths.''
\url{https://www.erdosproblems.com/65} (accessed 9 October 2026).
\bibitem{Dirac}
G. A. Dirac,
\emph{Some theorems on abstract graphs},
Proc. London Math. Soc. (3) \textbf{2} (1952), 69--81.
\url{https://doi.org/10.1112/plms/s3-2.1.69}.
\bibitem{Brandt}
S. Brandt, \emph{A sufficient condition for all short cycles},
Discrete Appl. Math. \textbf{79} (1997), 63--66.
\url{https://doi.org/10.1016/S0166-218X(97)00032-2}.
\bibitem{EG}
P. Erd\H{o}s and T. Gallai,
\emph{On maximal paths and circuits of graphs},
Acta Math. Acad. Sci. Hungar. \textbf{10} (1959), 337--356.
\url{https://doi.org/10.1007/BF02024498}.
\bibitem{MoonMoser}
J. Moon and L. Moser, \emph{On Hamiltonian bipartite graphs},
Israel J. Math. \textbf{1} (1963), 163--165.
\url{https://doi.org/10.1007/BF02759704}.
\end{thebibliography}
\end{document}
